% GENERATED FILE. Edit canonical lessons, then run npm run generate:books.
% canonical-source-hash: fnv1a64:9878d1f0a42bfe20
% canonical-lessons: PT-C134-nada, PT-C134-nenhum, PT-C134-ninguem, PT-C134-qual, PT-C134-nem

\chapter{Nothing, None, Nobody, Which One, Not Even}
\label{ch:pt-nothing-none-nobody-which-one-not-even}
\hlchaptermodality{134}

\begin{chapteropening}
\textbf{By the end of this chapter:} \emph{I can say that there is nothing, nobody and none, ask which one, and answer \textquotedblleft{}me neither\textquotedblright{}.}

Complete the last lesson of chapter 134: I can say that there is nothing, nobody and none, ask which one, and answer \textquotedblleft{}me neither\textquotedblright{}.
\end{chapteropening}

\section[nada]{nada — nothing}

\label{lesson:PT-C134-nada}



\begin{quote}
Before the new one: say the Portuguese for to tell, then the Portuguese for to walk.
\end{quote}

\subsection*{You'll want to know: nada}
\textbf{nada} — \textquotedblleft{}nothing\textquotedblright{}. With a verb, Portuguese keeps \textbf{não} in front as well: \textbf{Não vejo nada.} — \textquotedblleft{}I don't see anything.\textquotedblright{} It is the same word as in \textbf{de nada}: \textquotedblleft{}it's nothing\textquotedblright{}, so \textquotedblleft{}you're welcome\textquotedblright{}.

\begin{grammarlens}[title={What you've built}]
One more word for this chapter.
\end{grammarlens}

\subsection*{Your turn}
\begin{itemize}
\raggedright
  \item \emph{Say it:} \emph{nada}
  \item \emph{Say it:} \emph{nada}, once more
  \item \emph{Say it:} say \emph{nada}
  \item \emph{Recall:} say the Portuguese for to tell, then the Portuguese for to walk, then say \emph{nada} again
\end{itemize}

\begin{tcolorbox}[breakable,colback=teal!4,colframe=teal!35!black,title={Before you move on}]
What does nada mean? (\textquotedblleft{}Nothing\textquotedblright{}.) Say it once more.
\end{tcolorbox}
\section[nenhum / nenhuma]{nenhum / nenhuma — none, not a single}

\label{lesson:PT-C134-nenhum}



\begin{quote}
Before the new one: say the Portuguese for to walk, then the Portuguese for nothing.
\end{quote}

\subsection*{You'll want to know: nenhum / nenhuma}
\textbf{nenhum / nenhuma} — \textquotedblleft{}none, not a single\textquotedblright{}. It agrees with its noun, and it also keeps \textbf{não}: \textbf{Não há nenhum comboio ao sábado.} — \textquotedblleft{}There is no train at all on Saturdays.\textquotedblright{}

\begin{grammarlens}[title={What you've built}]
One more word for this chapter.
\end{grammarlens}

\subsection*{Your turn}
\begin{itemize}
\raggedright
  \item \emph{Say it:} \emph{nenhum / nenhuma}
  \item \emph{Say it:} \emph{nenhum / nenhuma}, once more
  \item \emph{Say it:} say \emph{nenhum / nenhuma}
  \item \emph{Recall:} say the Portuguese for to walk, then the Portuguese for nothing, then say \emph{nenhum / nenhuma} again
\end{itemize}

\begin{tcolorbox}[breakable,colback=teal!4,colframe=teal!35!black,title={Before you move on}]
What does nenhum / nenhuma mean? (\textquotedblleft{}None, not a single\textquotedblright{}.) Say it once more.
\end{tcolorbox}
\section[ninguém]{ninguém — nobody, no one}

\label{lesson:PT-C134-ninguem}



\begin{quote}
Before the new one: say the Portuguese for nothing, then the Portuguese for none.
\end{quote}

\subsection*{You'll want to know: ninguém}
\textbf{ninguém} — \textquotedblleft{}nobody, no one\textquotedblright{}. \textbf{Não vejo ninguém.} — \textquotedblleft{}I don't see anybody.\textquotedblright{} On its own it answers a question: \textbf{Ninguém.}

\begin{grammarlens}[title={What you've built}]
One more word for this chapter.
\end{grammarlens}

\subsection*{Your turn}
\begin{itemize}
\raggedright
  \item \emph{Say it:} \emph{ninguém}
  \item \emph{Say it:} \emph{ninguém}, once more
  \item \emph{Say it:} say \emph{ninguém}
  \item \emph{Recall:} say the Portuguese for nothing, then the Portuguese for none, then say \emph{ninguém} again
\end{itemize}

\begin{tcolorbox}[breakable,colback=teal!4,colframe=teal!35!black,title={Before you move on}]
What does ninguém mean? (\textquotedblleft{}Nobody, no one\textquotedblright{}.) Say it once more.
\end{tcolorbox}
\section[qual]{qual — which, which one}

\label{lesson:PT-C134-qual}



\begin{quote}
Before the new one: say the Portuguese for none, then the Portuguese for nobody.
\end{quote}

\subsection*{You'll want to know: qual}
\textbf{qual} — \textquotedblleft{}which, which one\textquotedblright{}. A question that picks from a set: \textbf{Qual é a tua mala?} — \textquotedblleft{}Which is your suitcase?\textquotedblright{} For more than one, \textbf{quais}: \textbf{Quais?} — \textquotedblleft{}Which ones?\textquotedblright{}

\begin{grammarlens}[title={What you've built}]
One more word for this chapter.
\end{grammarlens}

\subsection*{Your turn}
\begin{itemize}
\raggedright
  \item \emph{Say it:} \emph{qual}
  \item \emph{Say it:} \emph{qual}, once more
  \item \emph{Say it:} say \emph{qual}
  \item \emph{Recall:} say the Portuguese for none, then the Portuguese for nobody, then say \emph{qual} again
\end{itemize}

\begin{tcolorbox}[breakable,colback=teal!4,colframe=teal!35!black,title={Before you move on}]
What does qual mean? (\textquotedblleft{}Which, which one\textquotedblright{}.) Say it once more.
\end{tcolorbox}
\section[nem]{nem — nor, not even}

\label{lesson:PT-C134-nem}



\begin{quote}
Before the new one: say the Portuguese for nobody, then the Portuguese for which.
\end{quote}

\subsection*{You'll want to know: nem}
\textbf{nem} — \textquotedblleft{}nor, not even\textquotedblright{}. \textbf{nem ... nem} is \textquotedblleft{}neither ... nor\textquotedblright{}: \textbf{Não bebo nem café nem chá.} — \textquotedblleft{}I drink neither coffee nor tea.\textquotedblright{} After a no, \textbf{nem eu} means \textquotedblleft{}me neither\textquotedblright{}.

\begin{grammarlens}[title={What you've built}]
One more word for this chapter.
\end{grammarlens}

\subsection*{Your turn}
\begin{itemize}
\raggedright
  \item \emph{Say it:} \emph{nem}
  \item \emph{Say it:} \emph{nem}, once more
  \item \emph{Say it:} say \emph{nem}
  \item \emph{Recall:} say the Portuguese for nobody, then the Portuguese for which, then say \emph{nem} again
\end{itemize}

\begin{tcolorbox}[breakable,colback=teal!4,colframe=teal!35!black,title={Before you move on}]
What does nem mean? (\textquotedblleft{}Nor, not even\textquotedblright{}.) Say it once more.
\end{tcolorbox}
